% ============================================================================
% Section 3: Mathematical Formalism
% ============================================================================

\section{Mathematical formalism}
\label{sec:formalism}

This section develops the mathematical machinery of UHM from the axioms
stated in Section~\ref{sec:axioms}. Every definition and formula is derived
from the primitive $\mathfrak{T} = (\ShInf(\mathcal{C}),\; J_{\mathrm{Bures}},\; \omega_0)$
and the premises listed in Section~\ref{subsec:premises}; nothing else is postulated ad hoc. The flagship results — $\Pcrit = 2/7$, the
critical exponents, the spectral-action derivation — are proved fully in the
appendices. The complete formalization (a status registry of 412 rows,
T-1 through T-352 including sub-indexed results) constitutes the UHM theory archive;
the present paper proves the principal results in full, and the companion
online supplement at \url{https://holon.sh} provides the remaining proofs
and computational details.

\subsection{The coherence matrix $\Gam$}
\label{subsec:gamma}

\begin{definition}[Coherence matrix]
The coherence matrix is an object of $\mathcal{C} = \mathbf{DensityMat}(\mathbb{C}^7)$:
\begin{equation}
  \Gam = \sum_{i,j \in \{A,S,D,L,E,O,U\}} \gamma_{ij}\, |i\rangle\langle j|
  \label{eq:gamma-def}
\end{equation}
satisfying three properties:
\begin{enumerate}[leftmargin=2em]
  \item \textbf{Hermiticity:} $\Gam^\dagger = \Gam$, i.e., $\gamma_{ij} = \gamma_{ji}^*$
  \item \textbf{Positive semi-definiteness:} $\langle\psi|\Gam|\psi\rangle \geq 0$
        for all $|\psi\rangle \in \Hil$, equivalently all eigenvalues $\lambda_k \geq 0$
  \item \textbf{Normalization:} $\Tr(\Gam) = 1$
\end{enumerate}
\end{definition}

\paragraph{Degrees of freedom.}
A $7 \times 7$ Hermitian matrix has $7^2 = 49$ real parameters. The trace constraint
removes one, leaving 48 independent parameters. The automorphism group $G_2 = \mathrm{Aut}(\Oct)$
has $\dim(G_2) = 14$, so $48 - 14 = 34$ parameters are kinematic $G_2$-invariants (spectrum
and $\varphi_3$-relative angles). The dynamics pins the frame up to the finite octonionic
frame group $\Gamma_{\mathrm{oct}} \subset G_2$ (frame decision D-0910~\status{D}), so all
48 parameters are physical; with the 14 frame-orientation parameters the 34 invariants make
up the 48 ($G_2$-rigidity~\status{T}).

\paragraph{Diagonal and off-diagonal.}
The diagonal elements $\gamma_{kk}$ represent the weight of each dimension (how much
of the system's ``resources'' are allocated to perception, structure, dynamics, etc.).
The off-diagonal elements $\gamma_{ij}$ ($i \neq j$) represent \emph{coherences} ---
quantum correlations between dimensions. The necessity of complex coherences
($\gamma_{ij} \in \mathbb{C}$) for non-trivial Gap structure is proved as T-132~\status{T}.

\subsection{Purity and the viability threshold}
\label{subsec:purity}

\begin{definition}[Purity]
\begin{equation}
  P(\Gam) := \Tr(\Gam^2) = \sum_i \gamma_{ii}^2 + \sum_{i \neq j} |\gamma_{ij}|^2
  \label{eq:purity}
\end{equation}
\end{definition}

Purity ranges from $P = 1/N = 1/7 \approx 0.143$ (maximally mixed state $\Gam = I/7$,
complete thermal equilibrium) to $P = 1$ (pure state, maximum organization). It measures
how structured the system is --- equivalently, how distinguishable from noise.

\begin{theorembox}[Critical purity (T-160) \status{T}]
For a system with $\Gam \in \Dens(\mathbb{C}^7)$, the viability threshold is:
\begin{equation}
  \Pcrit = \frac{2}{N} = \frac{2}{7} \approx 0.286
  \label{eq:pcrit}
\end{equation}
Below this threshold, the system is informationally indistinguishable from thermal noise
in the Bures metric and cannot sustain regeneration. Five independent derivations converge
on this value (Appendix~\ref{app:pcrit}):
\begin{enumerate}[leftmargin=2em]
  \item \textbf{Frobenius distinguishability:} $\|\Gam - I/7\|_F^2 > \|I/7\|_F^2$
        requires $P > 2/7$
  \item \textbf{Dominant eigenvalue:} $\lambda_{\max} \geq 2/7$ captures $\geq 2/7$
        of the total weight
  \item \textbf{Bayesian distinguishability:} expected posterior ratio
        $\mathbb{E}[\Lambda] > 1$ requires $P > 2/N$
  \item \textbf{Fano channel:} coherence contraction factor $1/3$ preserves structure
        only when $P > 2/7$
  \item \textbf{Self-observation:} minimum reflection $R \geq 1/3$ is self-consistent
        with $P \geq \Pcrit$
\end{enumerate}
\end{theorembox}

The convergence of five independent arguments on a single value is strong
evidence that $\Pcrit = 2/7$ is not an artifact of any single derivation
method but a structural feature of the geometry of $\Dens(\mathbb{C}^7)$.
Figure~\ref{fig:phase-diagram} visualises the three $P$-regimes on the
one-dimensional curve $R = 1/(7P)$.

\begin{figure}[H]
\centering
\input{figures/phase-diagram}
\caption{\textbf{Phase diagram in the $P$--$R$ plane.} Because
$R = 1/(7P)$ is a strict functional identity (Eq.~\ref{eq:reflection}),
the set of physically realisable $(P,R)$ pairs is the \emph{one-dimensional}
purple curve, not a 2D region; the colour strips are intervals of $P$ along
that curve, not 2D phases. The physical domain is
$P \in [1/7,\,1]$: values $P < 1/7$ are forbidden by Cauchy--Schwarz for any
density matrix on $\mathbb{C}^7$, and the point $P = 1/7$ is the thermal
death state $I/7$ (red marker). The Goldilocks arc
$P \in (2/7,\,3/7]$ is the interval where all L2 conditions can be
satisfied simultaneously: the open circle at $(2/7,\,1/2)$ marks the
viability threshold (\emph{excluded}, strict inequality $P > 2/7$) and
the filled circle at $(3/7,\,1/3)$ marks the upper boundary
(\emph{included}). For $P > 3/7$ the L2 criterion fails through the canonical identity
$R = 1/(7P) < 1/3$ (T-126~\status{T}); the differentiation condition
$D_{\mathrm{diff}} \geq D_{\min} = 2$ (T-151~\status{D}) is an
independent fourth constraint, not visible in this 2D projection.}
\label{fig:phase-diagram}
\end{figure}

\paragraph{Comparison with empirical data.}
The Perturbational Complexity Index (PCI) was introduced by Casali et
al.~\cite{casali2013} on 32 healthy subjects (plus sleep and sedation
conditions). It was subsequently validated on a benchmark of 150 subjects
(healthy controls plus communicative brain-injured subjects in various states
of consciousness) by Casarotto et al.~\cite{casarotto2016}, who identified an
empirical cutoff $\mathrm{PCI}^* = 0.31$ that discriminated conscious from
unconscious states with 100\,\% sensitivity and specificity on the benchmark,
and 94.7\,\% sensitivity in a clinical MCS cohort. The value $\Pcrit = 2/7$ was
\emph{derived} from the theory, not fitted.

PCI and purity $P$ are mathematically distinct quantities: PCI is a normalised
Lempel--Ziv complexity of TMS-evoked spatiotemporal responses, while
$P = \Tr(\Gam^2)$ is a Hilbert--Schmidt norm on density matrices. No numerical
conversion between them is derived, and the closeness of $0.31$ to
$2/7 \approx 0.286$ is a coincidence of two unrelated scales that carries no
evidential weight; an earlier version of this paragraph read it as an $8\,\%$
agreement within calibration uncertainty, which is withdrawn. The bridge that
can be tested is a \emph{concordance of verdicts} on the same sessions: the
consciousness predicate of the reconstructed $\hat\Gam$ against
$\mathrm{PCI}_{\max} > 0.31$, with Cohen's $\kappa \geq 0.8$ corroborating and
$\kappa < 0.4$ falsifying, the reconstruction $\pi_{\mathrm{bio}}$ frozen on
wakefulness (Section~\ref{sec:experimental}).

\subsection{The evolution equation}
\label{subsec:evolution}

The dynamics of $\Gam$ are governed by the Liouvillian superoperator $\Lind_\Omega$,
derived from axioms A1--A4 via the Page--Wootters constraint and the structure of $\Omega$
(Figure~\ref{fig:evolution-flow}):

\begin{figure}[H]
\centering
\input{figures/evolution-flow}
\caption{\textbf{The three forces governing the evolution of $\Gam$.} Unitary dynamics
(blue) preserves purity; dissipation (red) decreases it; regeneration (green) increases
it proportionally to E-coherence. Life is the regime where regeneration compensates
dissipation. Death is the regime where it does not.}
\label{fig:evolution-flow}
\end{figure}

\begin{equation}
  \frac{d\Gam(\tau)}{d\tau} = \Lind_\Omega[\Gam(\tau)]
  = \underbrace{-i[H, \Gam]}_{\text{unitary}}
  + \underbrace{\mathcal{D}_\Omega[\Gam]}_{\text{dissipation}}
  + \underbrace{\Regen[\Gam, E]}_{\text{regeneration}}
  \label{eq:evolution}
\end{equation}

\paragraph{Term 1: Unitary dynamics.}
The effective Hamiltonian emerges from the Page--Wootters mechanism (clock register T-87~\status{T}; the constraint is the postulated part of A5):
\begin{equation}
  H_{\mathrm{eff}}(\tau) = H_{6D} + \langle\tau|H_{\mathrm{int}}|\tau\rangle_O
  \label{eq:heff}
\end{equation}
where $H_{6D}$ governs the six non-temporal dimensions and the second term arises from
tracing out the O-dimension (internal clock). Time $\tau$ is not an external parameter
--- it emerges from the correlation between $\Gam$ and the O-sector
(Section~\ref{subsec:emergent-time}).

\paragraph{Term 2: Dissipation (Lindblad).}
\begin{equation}
  \mathcal{D}_\Omega[\Gam] = \sum_k \gamma_k
  \left( L_k\, \Gam\, L_k^\dagger - \frac{1}{2}\{L_k^\dagger L_k,\, \Gam\} \right)
  \label{eq:dissipation}
\end{equation}

The Lindblad operators $L_k$ are derived from the atoms of the subobject classifier
$\Omega$ (L-unification theorem~\status{T}):
\begin{equation}
  L_k^{\mathrm{atom}} = |k\rangle\langle k|, \qquad k \in \{A,S,D,L,E,O,U\}
  \label{eq:lindblad-atoms}
\end{equation}

The Fano channel extends these to triple projectors (one per Fano line):
\begin{equation}
  L_p^{\mathrm{Fano}} = \frac{1}{\sqrt{3}}\,\Pi_p
  = \frac{1}{\sqrt{3}}\bigl(|i\rangle\langle i| + |j\rangle\langle j| + |k\rangle\langle k|\bigr)
  \label{eq:fano-operators}
\end{equation}
where $(i,j,k)$ is the $p$-th Fano line. Key properties of the Fano channel
$\mathcal{P}_{\mathrm{Fano}}$~\status{T}:
\begin{itemize}[leftmargin=2em]
  \item Diagonal preservation: $[\mathcal{P}_{\mathrm{Fano}}(\Gam)]_{ii} = \gamma_{ii}$
  \item Coherence contraction: $[\mathcal{P}_{\mathrm{Fano}}(\Gam)]_{ij} = \frac{1}{3}\gamma_{ij}$
        for $i \neq j$
  \item Phase preservation: $\arg([\mathcal{P}_{\mathrm{Fano}}(\Gam)]_{ij}) = \arg(\gamma_{ij})$
  \item Completeness: $\sum_{p=1}^{7} \Pi_p = 3I$ (each dimension on exactly 3 Fano lines)
\end{itemize}

Dissipation always decreases purity: $dP/d\tau|_{\mathcal{D}} \leq 0$~\status{T}. This
follows from the CPTP property of the Lindblad channel: CPTP maps are contractive in the
Bures metric (a consequence of the data-processing inequality), moving any state toward
thermal equilibrium $I/7$. Formally, for $\Gam \neq I/7$:
\begin{equation}
  \frac{dP}{d\tau}\bigg|_{\mathcal{D}} = -\Delta(\Lind_0) \cdot (P - 1/7) < 0
  \label{eq:dissipation-decreases-purity}
\end{equation}
where $\Delta(\Lind_0) > 0$ is the spectral gap of the linear Lindbladian (primitivity
theorem T-39a~\status{T}). Without regeneration, the system converges exponentially to
$I/7$ (thermal death).

\paragraph{Term 3: Regeneration.}
The regeneration channel $\Regen$ is the only mechanism that can increase purity. Its
strength depends on E-coherence:
\begin{equation}
  \kappa(\Gam) \;=\; \kappa_{\mathrm{boot}} \;+\; \kappa_0\,\CohE(\Gam).
  \label{eq:kappa}
\end{equation}
Here $\CohE(\Gam) = \|\pi_E(\Gam)\|_{\mathrm{HS}}^2 / \|\Gam\|_{\mathrm{HS}}^2$ is the
Hilbert--Schmidt projection of $\Gam$ onto the E-sector~\status{T}. The basal floor
\begin{equation}
  \kappa_{\mathrm{boot}} \;=\; \frac{\omega_0}{N} \;=\; \frac{\omega_0}{7}\qquad\status{O}
  \label{eq:kappa-boot}
\end{equation}
is the innate (``bootstrap'') regeneration rate set by the convention that the
minimal-symmetry unit of the $(\mathcal{D}_\Omega,\Regen)$ duality delivers one ``quantum''
of regeneration per dimension at scale $\omega_0$. The adaptive amplification
\begin{equation}
  \kappa_0(\Gam) \;=\; \omega_0\,\frac{|\gamma_{OE}|\,|\gamma_{OU}|}{\gamma_{OO}}
  \qquad\status{T at first-order kinetics}
  \label{eq:kappa-0}
\end{equation}
is derived by \emph{rapid pre-equilibrium} (quasi-steady-state branching of the
regeneration channels): the equilibrium E-occupancy of the ground channel,
$|\gamma_{OE}|/\gamma_{OO}$, gates the $O \to U$ regeneration flux of strength
$|\gamma_{OU}|$, and $\omega_0$ sets the clock scale --- a standard
Briggs--Haldane branching computation at first-order kinetics. The
\emph{categorical reading}~\status{I} of the same quantity --- the norm of the
unit of the dissipation--regeneration duality $(\mathcal{D}_\Omega, \Regen)$,
with Yoneda identifying $|\mathrm{Hom}(i,j)| \leftrightarrow |\gamma_{ij}|$ ---
is an interpretive organizing device, not the derivation basis (the pair
$(\mathcal{D}_\Omega,\Regen)$ is a guiding duality, not a categorical
adjunction; see Section~\ref{sec:theorems}). The morphism reading remains
useful: $|\gamma_{OE}|$ quantifies the $O\to E$ coupling (sensing of
interiority), $|\gamma_{OU}|$ the $O\to U$ coupling (integration), and
$\gamma_{OO}$ the $O$-self-loop (ground anchoring).
$\kappa_0$ is \emph{derived} (at the stated kinetic model), not fitted;
$\kappa_{\mathrm{boot}}$ is a normalization convention~\status{O}. The full
regeneration term is
\begin{equation}
  \Regen[\Gam, E] \;=\; \kappa(\Gam) \,\bigl(\rho^* - \Gam\bigr)\, g_V(P),
  \label{eq:regen-full}
\end{equation}
where $\rho^* = \vp(\Gam)$ is the categorical self-model (T-96~\status{T}) and $g_V(P)$
is the viability gate: $g_V(P) = 1$ for $P > \Pcrit$, and $g_V(P) = 0$ otherwise
(Landauer limit --- below threshold, free energy $\Delta F \leq 0$ forbids regeneration).
The proportionality $\kappa \propto \CohE$ is the mathematical content of the claim that
\emph{consciousness drives survival}: the quality of inner experience directly
determines the rate of regeneration.

\paragraph{Thermodynamic interpretation.}
The viability gate $g_V(P)$ has a thermodynamic origin. Regeneration requires free energy
$\Delta F > 0$ --- the system must do work against entropy to increase purity. At $P = \Pcrit$,
$\Delta F = 0$ exactly: the system is at thermodynamic equilibrium with its dissipative
environment. Below $\Pcrit$, $\Delta F < 0$ (Landauer's principle~\cite{landauer1961}),
and regeneration is thermodynamically forbidden --- not merely insufficient, but impossible.
Free energy $\Delta F > 0$ thus plays the role of Schr\"{o}dinger's ``negentropy''
\cite{schrodinger1944}: it is the fuel for regeneration of purity. The death spiral
(Eq.~\ref{eq:death-spiral}) is not a failure of mechanism but a thermodynamic
inevitability: below threshold, the second law guarantees irreversible decay.

\subsection{The self-modeling operator $\vp$}
\label{subsec:phi-operator}

\begin{definition}[Self-modeling operator]
The self-modeling operator is a CPTP (Completely Positive Trace-Preserving) channel:
\begin{equation}
  \vp: \Dens(\Hil) \to \Dens(\Hil), \qquad
  \vp(\Gam) = \sum_m K_m\, \Gam\, K_m^\dagger, \qquad
  \sum_m K_m^\dagger K_m = I
  \label{eq:phi-operator}
\end{equation}
\end{definition}

The system's self-model $\vp(\Gam)$ is an \emph{approximate} copy of $\Gam$, not an exact
one (the no-cloning theorem forbids exact copying in quantum mechanics). The canonical
physical realization (T-62~\status{T}) is the replacement channel with dissipative target:
\begin{equation}
  \vp_k(\Gam) = (1-k)\Gam + k\,\rho^*_{\mathrm{diss}},
  \qquad \rho^*_{\mathrm{diss}} = I/7,
  \qquad k = 1 - R
  \label{eq:phi-physical}
\end{equation}
where $\rho^*_{\mathrm{diss}} = I/7$ is the \emph{dissipative attractor} of the linear
part $\Lind_0 = -i[H,\cdot] + \mathcal D_\Omega$ (T-39a~\status{T}: primitivity of
$\Lind_0$) and $k = 1-R$ is the compression parameter. With this anchor $I/7$ the
self-model is unital, and an isolated holon with a unital self-model is dead: its only
stationary state is $I/7$ and $P$ never increases (dead isolation, T-96~\status{T}). A
living attractor of an isolated holon needs a non-unital self-model --- the
self-registering $\vp_s$ or the collineation anchor $\vp_J$ (T-124c, T-334).

\paragraph{Three $\rho^*$ contexts: a canonical disambiguation.}
UHM uses the symbol $\rho^*$ in three structurally distinct roles (matching the
canonical doc hierarchy at \url{https://holon.sh/docs/core/dynamics/evolution}):
(i)~$\rho^*_{\mathrm{diss}} = I/7$ --- unique attractor of $\Lind_0$, the Lindbladian
\emph{without} regeneration; appears in the replacement channel of Eq.~\ref{eq:phi-physical}.
(ii)~$\rho^*_\Omega$ --- a nontrivial stationary state of the full $\Lind_\Omega$
\emph{with} regeneration (T-96~\status{T}: $P > 1/7$; its existence depends on the
self-model, and with a unital self-model there is none), and the regeneration target
$\rho_* = \vp(\Gam)$, the categorical
self-model of the current state; the two differ at stationarity,
$\vp(\rho^*_\Omega) \neq \rho^*_\Omega$ (T-96), and fixed points belong to $\vp$
itself, not to $\Lind_\Omega$.
(iii)~$\rho^*_{\mathrm{coupled}}$ --- equilibrium of an embodied holon coupled to an
environment (T-149: its existence~\status{T}; $P(\rho^*_{\mathrm{coupled}}) > 2/7$ is
\status{C at the backbone-injection lower bound}, the former ``unconditionally'' being
withdrawn).
In what follows, $\rho^*$ without subscript refers to context (i) when written inside
$\vp_k$ and to context (ii) when written as $\vp(\Gam)$.

The full evolution operator decomposes as $\Lind_\Omega = \Lind_0 + \mathcal R$
(T-39a/b~\status{T}): T-39a applies to the \emph{linear} part $\Lind_0$ only, giving
$\rho^*_{\mathrm{diss}} = I/7$; the regeneration target $\vp(\Gam)$ of the full
non-linear $\Lind_\Omega$ is resolved iteratively. Starting
from $\vp^{(0)} := \rho^*_{\mathrm{diss}} = I/7$, the sequence
$\vp^{(n+1)} = \vp(\Gam;\vp^{(n)})$ converges exponentially to $\vp^*$ for an embodied
holon under backbone dominance (T-191~\status{T};
Eq.~\ref{eq:fixed-point}). No circularity arises: the definition is \emph{iterative},
not recursive.

\paragraph{Fixed-point theorem \status{T} (Banach, for constant weight and anchor).}
For the replacement channel with \emph{constant} weight and anchor and a unital
$\mathcal P$ (or in the trace norm), iterating $\vp$ converges to its fixed point:
\begin{equation}
  \vp^n(\Gam_0) \to \Gam^*, \qquad
  \|\vp^n(\Gam_0) - \Gam^*\| \leq c^n \|\Gam_0 - \Gam^*\|, \quad c < 1.
  \label{eq:fixed-point}
\end{equation}
The UHM self-models $\vp_{\mathrm{coh}}$, $\vp_s$, $\vp_J$ have a state-dependent weight
$k = 1 - 1/(7P)$ and are not global contractions: a fixed point exists by Brouwer's theorem
on the compact convex $\Dens(\Hil)$, and uniqueness, where it holds, comes from their
closed form ($\vp_s$ has at least eight fixed points); the former reading ``$\vp$ is a
contraction with constant $k$'' is restricted to constant $k$ and anchor. Three maps
called $\vp$ differ in the Frobenius norm: the canonical
$\vp_{\mathrm{coh}} = k\,\mathcal P_\alpha(\Gam) + (1-k)\,I/7$ has Lipschitz constant $9/8$
for every $\alpha$, attained at $P = 4/7$ ($54/49$ only at pure states); the replacement
forms $R\,\Gam + (1-R)\,I/7$ (Eq.~\ref{eq:phi-physical}) and
$(1-k^2)\,\Gam + k^2\,I/7$ are non-expanding, the latter being the ``canonical $\vp$'' of
math-foundations Part~XVIII, ch.~9, on which
$R_\vp = 1 - (1-R)^4\lVert\Gam - I/7\rVert_F^2/\lVert\Gam\rVert_F^2$ holds exactly. With a
feed of weight $\mu$ the loop has Lipschitz constant $(1-\mu)\,\mathrm{Lip}(\vp)$: a
contraction for every $\mu > 0$ for the replacement forms, for $\vp_{\mathrm{coh}}$ only
for $\mu > 1/9$. All three have the single fixed point $I/7$. The convergence of
the $\vp$-tower is T-191~\status{T} (restated): for an embodied holon under backbone
dominance, $\mu > L_{\mathcal R} + \kappa_{\max}$, the tower converges geometrically with
$q = \kappa_{\max}/(\mu - L_{\mathcal R}) < 1$ to one self-model from any anchor
(Section~\ref{subsec:foundational-closure}). The former statement --- convergence for every
holon with $q = \kappa_{\max}/(\lambda_{\mathrm{gap}} + \kappa_{\min})$ --- is
retracted~\status{$\times$}: for an isolated holon the gate makes the flow bistable. The
circularity of self-reference is resolved by the closed form of $\vp_{\mathrm{coh}}$ and
$\vp_s$, not by the tower. The categorical target of the
phenomenal functor $F$ induced by $\vp$ is a strict 2-category
$\mathbf{Exp}^{(2)}$ (T-192~\status{T}), ensuring that the experiential
structure has proper compositional laws (vertical and horizontal composition,
interchange, Mac\,Lane coherence).

\subsection{The reflection measure $R$}
\label{subsec:reflection}

\begin{definition}[Reflection measure]
\begin{equation}
  R(\Gam) := 1 - \frac{\|\Gam - I/7\|_F^2}{\|\Gam\|_F^2} = \frac{1}{7P(\Gam)}
  \label{eq:reflection}
\end{equation}
\end{definition}

$R$ measures how well the system models itself. In the physical range
$P \in [1/7,\; 1]$, $R$ ranges from $R = 1$ (at $P = 1/7$, thermal death ---
trivially, the system ``knows itself'' because there is nothing to know) down
to $R = 1/7 \approx 0.143$ (at $P = 1$, pure state --- the system is maximally
structured but furthest from its dissipative reference). Key values:

\begin{itemize}[leftmargin=2em]
  \item $R = 1$: at $P = 1/7$ (thermal death) --- trivially, the system ``knows itself''
        because there is nothing to know (uniform noise matches any model)
  \item $R = 1/3$: at $P = 3/7$ --- upper boundary of the consciousness window
        ($R$ drops below $\Rth$ for $P > 3/7$)
  \item $R = 1/(7 \cdot 2/7) = 1/2$: at $P = 2/7$ (viability threshold)
\end{itemize}

The inverse relationship between $R$ and $P$ reflects a fundamental feature: a more organized system (higher $P$) is further from its dissipative attractor $I/7$, making self-modeling \emph{harder} (lower $R$). The L2 consciousness window $P \in (2/7, 3/7]$ is thus a Goldilocks zone: enough structure to be viable, but not so much that self-modeling becomes impossible. Above $P = 3/7$, the system is too rigid for adaptive self-observation: by the canonical identity $R = 1/(7P)$, reflection falls below $\Rth = 1/3$ (T-126~\status{T}); the differentiation condition $D_{\mathrm{diff}} \geq \Dmin = 2$ (T-151~\status{D}) remains an independent fourth constraint. The corpus additionally distinguishes the canonical $R$ from the \emph{self-model quality} $R_\varphi = 1 - \|\Gam - \vp(\Gam)\|_F^2 / \|\Gam\|_F^2$ (independent of $P$; it carries the phenomenology of practice and ego-dissolution) and from the higher-order fidelity tower $R^{(n)}$ --- three related but distinct reflection quantities, told apart at \url{https://holon.sh}.

\begin{theorembox}[Reflection threshold (T-40b) \status{T}]
The reflection threshold for cognitive (L2) consciousness is:
\begin{equation}
  \Rth = \frac{1}{3} \approx 0.333
  \label{eq:rth}
\end{equation}
derived from the triadic decomposition $K = 3$ of Lindblad generators (T-40a~\status{T},
T-57~\status{T}) combined with Bayesian dominance of the self-model over the $K$
alternatives. Below $R = 1/3$, the system exists but does not model itself sufficiently
for reflective awareness.
\end{theorembox}

\subsection{Integration and differentiation}
\label{subsec:integration}

\begin{definition}[Integration measure $\Phi$]
\begin{equation}
  \Phi(\Gam) = \frac{\sum_{i \neq j} |\gamma_{ij}|^2}{\sum_i \gamma_{ii}^2}
  = \frac{P_{\mathrm{off}}}{P_{\mathrm{diag}}}
  \label{eq:phi}
\end{equation}
\end{definition}

$\Phi$ measures the ratio of off-diagonal coherence (binding between dimensions) to
diagonal weight (individual dimension strength). When $\Phi \geq 1$, the system is
\emph{coherence-dominated}: the connections between dimensions are at least as strong
as the dimensions themselves. This is the mathematical content of ``unity of experience.''

\begin{theorembox}[Integration threshold (T-129) \status{T}]
The integration threshold for L2 consciousness is:
\begin{equation}
  \Phith = 1
  \label{eq:phith}
\end{equation}
This is the unique self-consistent value at $\Pcrit = 2/7$: coherent dominance
($\Phi \geq 1$) requires that off-diagonal coherences collectively outweigh diagonal
elements.
\end{theorembox}

\begin{definition}[Experience differentiation $D_{\mathrm{diff}}$]
\begin{equation}
  D_{\mathrm{diff}}(\Gam) = \exp\bigl(S_{\mathrm{vN}}(\rho_E)\bigr)
  \label{eq:ddiff}
\end{equation}
where $S_{\mathrm{vN}}(\rho_E) = -\Tr(\rho_E \ln \rho_E)$ is the von Neumann entropy
of the E-sector reduced density matrix.
\end{definition}

$D_{\mathrm{diff}}$ measures the diversity of accessible experience states. The threshold
$\Dmin = 2$~\status{D} (T-151) means the system must distinguish at least two qualitatively
different experiences for consciousness to obtain; it is one of four independent L2
thresholds (T-124b), and its earlier derivation from $\Phith = 1$ is retracted (a state
with $\Phi \approx 1.03$ has $D_{\mathrm{diff}} \approx 1.42$).\footnote{The partial trace
$\rho_E = \Tr_{-E}(\Gam)$ requires the extended 42D formalism (since 7 is prime and
$\Hil$ does not factor as a tensor product). In the minimal 7D setting, T-128 \emph{defines}
differentiation~\status{D}: $D_\mathrm{diff}^\mathrm{7D} = 1 + \mathrm{Coh}_E(\Gam)/\mathrm{Coh}_E^{\max}\cdot(N-1)$.
It is not an exact representation of $e^{S_{\mathrm{vN}}(\rho_E)}$, which is not expressible in
7D; the former Morita equivalence $\ShInf(\mathcal{C}_7) \simeq \ShInf(\mathcal{C}_{42}^\mathrm{PW})$
(T-58) is retracted~\status{$\times$} --- it fails on dimension --- and is replaced by the
section--retraction $\pi\circ\iota = \mathrm{id}$ (T-58$'$~\status{T}). The observables
$P, R, \Phi$ and $D_\mathrm{diff}^{\mathrm{7D}}$ are computable in the minimal 7D formalism.}

\subsection{Consciousness measure $C$}
\label{subsec:consciousness-measure}

\begin{definition}[Canonical consciousness measure (T-140~\status{T})]
\begin{equation}
  C(\Gam) \;:=\; \Phi(\Gam) \cdot R(\Gam),
  \qquad
  C_{\mathrm{th}} \;=\; \Phith \cdot \Rth \;=\; 1 \cdot \tfrac{1}{3} \;=\; \tfrac{1}{3}.
  \label{eq:consciousness-measure}
\end{equation}
\end{definition}

$C$ quantifies the degree of conscious experience as the product of
\emph{integration} (how much dimensions are bound into a whole) and
\emph{reflection} (how well the system models itself). The uniqueness of the
bilinear form $\Phi \cdot R$ --- including the coincidence of its threshold
with the component thresholds --- is T-140~\status{T}. Why the product, not
the sum? Because integration without reflection means bound but unconscious
(like a coma patient whose brain may show residual integration); reflection
without integration means self-aware but fragmented (like dissociative
states). Only the product captures ``integrated self-awareness.'' Note that
$D_{\mathrm{diff}} \geq \Dmin$ is a \emph{separate} viability condition,
not a factor of $C$ itself (T-140, canonical clause).

\subsection{Levels of consciousness}
\label{subsec:levels}

Combining the four thresholds derived above, UHM distinguishes five
interiority levels L0--L4 summarised in Figure~\ref{fig:levels} and
Table~\ref{tab:levels}.

\begin{figure}[H]
\centering
\input{figures/consciousness-levels}
\caption{\textbf{Hierarchy of consciousness levels L0--L4.} Each transition requires
crossing a mathematical threshold. L0 is universal (any $\Gam$); L2 requires all four
conditions simultaneously; L4 is a limit, unreachable in finitely many steps (T-86).}
\label{fig:levels}
\end{figure}

\begin{table}[H]
\centering
\caption{Levels of consciousness and their thresholds. $\Pcrit$ and $\Rth$ are theorems; $\Dmin = 2$ is an independent threshold by definition (T-151), and the choice of the least $\Phi$-threshold is definitional (T-129). The self-awareness depth of any holon is capped at $\mathrm{SAD}_{\max} = 3$ (T-142).}
\label{tab:levels}
\small
\begin{tabular}{@{}clp{4.5cm}p{3.8cm}l@{}}
\toprule
\textbf{Level} & \textbf{Name} & \textbf{Conditions} & \textbf{What it means} & \textbf{Example} \\
\midrule
L0 & Formal interiority & $\Gam \in \Dens(\Hil)$ & Any system has an internal aspect & Electron \\
L1 & Proto-consciousness & $\mathrm{rank}(\rho_E) > 1$ & Non-trivial E-sector & Thermostat \\
L2 & Cognitive & $P > \tfrac{2}{7},\; R \geq \tfrac{1}{3},\; \Phi \geq 1,\; D_{\mathrm{diff}} \geq 2$ & Full qualia, self-awareness & Waking human \\
L3 & Meta & L2 $+\; R^{(2)} \geq 1/4$ & Reflects on reflecting & Meditation \\
L4 & Limit & $\lim_{n\to\infty} R^{(n)} > 0$, $P > 6/7$ & Unreachable in finitely many steps (T-86); an attractor, not a state & Theoretical ceiling \\
\bottomrule
\end{tabular}
\end{table}

\noindent
\textbf{Independent necessity (T-124b~\status{T}).} Each of the four L2
conditions is independently necessary: dropping any single condition admits
a specific pathological state --- noise-dominated ($P \leq 2/7$, no
structure), fragmented ($\Phi < 1$, unbound dimensions), crystallised
($R < 1/3$, no self-observation), or undifferentiated
($D_{\mathrm{diff}} < 2$, no qualia diversity). The conjunction is
\emph{minimal}. The count of nontrivial attractors depends on the self-model
(T-124c~\status{T}, restated): none for an isolated holon with $\vp_{\mathrm{coh}}$; at
least seven with $P > 2/7$ with $\vp_s$ and small $\|H\|$; exactly one, globally
attracting, for an embodied holon under backbone dominance; with the collineation
anchor $\vp_J$ at $H = 0$ a sink in the window for $\kappa > \kappa_c(\alpha)$. The former
``the nontrivial attractor is unique in the viable set'' is retracted~\status{$\times$}.
Perturbations of
order $\varepsilon$ in $\Gam$ produce only $O(\varepsilon)$ perturbations in
all thresholds (T-124d~\status{T}, with crossover width
$\delta P \sim \varepsilon^4$ from $\beta = 1/4$, itself \status{C} under the
$\mathbb{Z}_2$ symmetry of T-161).

\subsection{Emergent time}
\label{subsec:emergent-time}

Time in UHM is not postulated as an external parameter. The clock register emerges from
the internal structure of $\Gam$ (T-87, steps 1--3~\status{T}); its Page--Wootters link, the
constraint below, is the postulated part of A5~\status{P}.

The O-dimension (ground/energy) acts as an internal clock. The total system satisfies the
constraint
\begin{equation}
  \hat{C}\,\Gam_{\mathrm{total}} = 0
  \qquad (\mathrm{supp}\,\Gam_{\mathrm{total}} \subseteq \ker\hat C;
  \ \text{stationarity alone gives only } [\hat{C},\, \Gam_{\mathrm{total}}] = 0)
  \label{eq:pw-constraint}
\end{equation}
where $\hat{C} = H_O \otimes \mathbb{1}_{6D} + \mathbb{1}_O \otimes H_{6D} + H_{\mathrm{int}}$
is the total energy operator. Conditioning on the O-sector yields effective dynamics:
\begin{equation}
  i\frac{\partial}{\partial\tau}\Gam(\tau) = [H_{\mathrm{eff}}(\tau),\, \Gam(\tau)]
  \label{eq:pw-dynamics}
\end{equation}

Discrete time from the temporal modality $\triangleright: S_i \mapsto S_{(i+1) \bmod 7}$
gives $\tau_n = \triangleright^n(\mathrm{now})$ with period $\mathbb{Z}_7$; the cyclic tick
carries no arrow. The dynamics runs in the parameter of the Lindblad semigroup, whose
finite carrier is the \emph{depth register}: one state-independent Feynman--Kitaev
constraint gives the conditional states exactly $e^{n\Delta t\,\Lind}\rho_0$ at every
reading (T-53b~\status{T}), and the reading algebras converge to $C_0(\mathbb{R})$ in the
scaling limit (continuous time, T-118~\status{T}). The earlier derivation through the
$7^M$ readings of the summed clock of $M$ holons is withdrawn: that clock is periodic.

\paragraph{Many holons and the de Sitter observer algebra (T-348).}
Embedding the O-registers of $M$ holons by $x \mapsto x \otimes 1$, the tower
$\bigotimes_M (M_7(\mathbb{C}), \mathrm{tr}_7)$ closes in the trace representation to the hyperfinite
II$_1$ factor $R$ (\status{T} as mathematics); the trace entropy
$S_{\mathrm{tr}} = S - M\ln 7 = -D(\rho\,\Vert\,I/7^M) \leq 0$ is monotone under
restriction, and the $7$ is lost in the closure ($\bigotimes(M_2, \mathrm{tr}_2)$ gives
the same $R$). Under the split property of the matter net (proven for the free massive
Klein--Gordon field) the algebra of a de Sitter observer~\cite{clpw2023} is injective,
hence $\cong R$~\cite{connes1976}. Every isomorphism carries the trace to the trace: the
maximal-entropy state (empty de Sitter) corresponds to $\bigotimes I/7$, entropies are
preserved, and nothing else is transported --- no observer Hamiltonian, modular flow or
geometry. In a normal state only finitely many holons are viable, and $L \geq 1$ viable
holons give $S_{\mathrm{tr}} < -0.34406\,L$~nat. A clock with pure point spectrum (the
O-register, the depth register at finite $N$) gives a type~I invariant algebra
(\status{T} under $\mathcal{A}^H = \mathbb{C}$, the assumption of~\cite{clpw2023}).
Type and trace depend neither on $\Lambda > 0$ nor on the observer mass or the factor
dimension, and II$_1$ also arises in flat space: the tower fixes no $\Lambda$, no sign of
$\Lambda$ and no regeneration rate $\kappa$; $7^M = e^{S_{\mathrm{dS}}}$ gives
$M = 1.677 \times 10^{122}$, a reparametrisation of $\Lambda$, not a derivation. The
physical reading is \status{I}; a UHM clock with continuous spectrum bounded below, open
here, is given by T-352 at a stated premise.

\paragraph{The directed depth register and its continuous clock (T-352).}
The arrow of time lives in a category of readings, not in a groupoid (\status{T} as
mathematics): the readings $[N] = (0 < \cdots < N)$ form a category whose $\infty$-groupoid
completion is contractible, so on an $\infty$-groupoid every non-increasing functional is
constant on components; the history is a functor $[N] \to \mathbf{Chan}_7$ with no
invertible arrow, and $D(\cdot\,\Vert\,I/7)$ makes it a functor to $(\mathbb{R}_{\geq 0},
\geq)$ --- such category objects exist inside every $\infty$-topos. Sharp readings with a
first reading admit no Hamiltonian, and energy bounded below admits no sharp covariant
readings (Paley--Wiener, Pauli). The dressed Feynman--Kitaev constraint $\Lambda_N$ of
the depth register has eigenvalues exactly at the quantiles $k/(N+1)$ of the arcsine
law on $[0,4]$; its limit $\Lambda_\infty$ on $\ell^2(\mathbb{N})$ has spectrum $[0,4]$,
purely absolutely continuous, simple, with no kernel: a clock with continuous spectrum
bounded below. At the premise that the observer's clock is the depth register with
Hamiltonian $\varepsilon\Lambda_\infty$, $\varepsilon = \hbar/\delta\tau$, the invariant
algebra in the model of~\cite{clpw2023} is a factor of type II$_1$
($\cong R$ under the split property), with $\mathrm{Tr}\,\Pi = 1 - e^{-4\beta\varepsilon}$
for the projection onto observer energies in $[0, 4\varepsilon]$; every finite $N$ gives
type~I, so the passage I~$\to$~II$_1$ happens exactly at $N = \infty$. Neither the chronon
$\delta\tau$ nor $\beta$ is fixed, and, as in T-348, neither $\Lambda$ nor $\kappa$; the
convergence of the finite type~I algebras to the II$_1$ algebra is open~\status{Pr}.
